\subsection{Characterizations of integrable numerical functions}

PROPOSITION 5. --- For a numerical function \(f \geqslant 0\) (finite or not), lower semi-continuous on X, to be integrable, it is necessary and sufficient that \(\int^{*} f d|\mu| < +\infty\) .

It all comes down to proving that the condition is sufficient. The definition of \(|\mu|^*(f)\) (§1, No.~1, Def. 1) proves that, for every \(\varepsilon > 0\) , there exists a continuous function \(g \geqslant 0\) , with compact support, such that \(g \leqslant f\) and \(|\mu|^*(f) \leqslant |\mu|(g) + \varepsilon\) . But \(f - g\) is lower semi-continuous and \(\geqslant 0\) , therefore (§1, No.~1, Th. 2)

\[
| \mu | ^ {*} (f) = | \mu | (g) + | \mu | ^ {*} (f - g),
\]

in other words \(\mathrm{N}_{1}(f-g)=|\mu|^{*}(f-g)=|\mu|^{*}(f)-|\mu|(g)\leqslant\varepsilon\) , which proves that f is integrable (§3, No.~4, Prop. 7).

COROLLARY 1. --- For a finite numerical function \(f \geqslant 0\) , upper semi-continuous on X, to be integrable, it is necessary and sufficient that \(\int^{*} f d|\mu| < +\infty\) .

For, if \(|\mu|^*(f) < +\infty\) , then there exists a lower semi-continuous function \(h\) such that \(f \leqslant h\) and \(|\mu|^*(h) < +\infty\) ; \(h - f\) is everywhere defined and lower semi-continuous, and \(|\mu|^*(h - f) \leqslant |\mu|^*(h) < +\infty\) ; therefore \(h - f\) is integrable, and since \(f(x) = h(x) - (h(x) - f(x))\) almost everywhere, \(f\) is integrable.

COROLLARY 2. --- Let H be a nonempty set, directed for the relation \(\leqslant (\text{resp. } \geqslant)\) , of lower (resp. upper) semi-continuous and integrable numerical functions; if

\[
\sup _ {f \in \mathrm{H}} \int f d | \mu | <   + \infty \quad (\mathrm{resp.} \inf _ {f \in \mathrm{H}} \int f d | \mu | > - \infty),
\]

then the upper (resp. lower) envelope g of H is integrable,

\[
\int g d \mu = \lim _ {f \in \mathrm{H}} \int f d \mu ,
\]

and \(\int g d|\mu| = \sup_{f\in \mathrm{H}}\int f d|\mu|\) (resp. \(\int g d|\mu| = \inf_{f\in \mathrm{H}}\int f d|\mu|\) ).

We may limit ourselves to the case of lower semi-continuous functions; the functions \(f^{+}\) (resp. \(f^{-}\) ), as \(f\) runs over \(\mathbf{H}\) , then form a directed set for \(\leqslant (\text{resp. } \geqslant)\) of lower (resp. upper) semi-continuous functions \(\geqslant 0\) ; the upper (resp. lower) envelope of the \(f^{+}\) (resp. \(f^{-}\) ), for \(f \in \mathbf{H}\) , is equal to \(g^{+}\) (resp. \(g^{-}\) ). On the other hand, one can replace \(\mathbf{H}\) by one of its sections (which is cofinal to it), consisting of the \(f \in \mathbf{H}\) that are \(\geqslant f_{0}\) , for some function \(f_{0} \in \mathbf{H}\) ; then \(\int f^{+} d|\mu| \leqslant \int f d|\mu| + \int f_{0}^{-} d|\mu|\) ; we thus see that we are reduced to proving the two assertions of the corollary when \(\mathbf{H}\) consists of positive functions. If \(\mathbf{H}\) is directed for \(\leqslant\) and consists of lower semi-continuous functions \(\geqslant 0\) , then we know (§1, No.~1, Th. 1) that

\[
| \mu | ^ {*} (g) = \sup _ {f \in \mathrm{H}} | \mu | ^ {*} (f) = \sup _ {f \in \mathrm{H}} \int f d | \mu | <   + \infty ,
\]

therefore g, which is lower semi-continuous, is integrable by Prop. 5; we have \(\int g d|\mu| = \lim_{f \in H} \int f d|\mu|\) and, since \(f \leqslant g\) , \(\lim_{f \in H} N_1(g - f) = 0\) , which shows that f converges in mean to g with respect to H, and thus proves the corollary in this case. If H is directed for \(\geqslant\) and consists of upper semi-continuous integrable functions f such that \(0 \leqslant f \leqslant f_1\) with \(f_1 \in H\) , then there exists a lower semi-continuous integrable function h such that \(f_1 \leqslant h\) ; we may write \(f = h - f'\) , where \(f'(x) = h(x) - f(x)\) when \(f(x) < +\infty\) , and \(f'(x) = 0\) otherwise. It is clear that the \(f'\) form a directed set, for \(\leqslant\) , of lower semi-continuous integrable functions \(\geqslant 0\) , with

\[
\int f ^ {\prime} d | \mu | \leqslant \int h d | \mu | <   + \infty ;
\]

we can apply to them what has been proved above; if \(g'\) is the upper envelope of the \(f'\) , then h and \(g'\) are finite almost everywhere, therefore \(h - g'\) is defined almost everywhere and is equal to g almost everywhere; from this, the conclusions of the corollary follow at once in this case.

COROLLARY 3. --- Let \(f\) be a bounded numerical function, upper semi-continuous on \(X\) and with compact support. Then, the mapping \(\mu \mapsto \int f d\mu\) is upper semi-continuous on \(\mathcal{M}_{+}(X)\) for the vague topology.

If \(h\) is a function in \(\mathcal{H}_{+}(\mathrm{X})\) such that \(|f| \leqslant h\) (Ch. III, §1, No.~2, Lemma 1) then \(0 \leqslant f + h \leqslant 2h\) , and since \(f + h\) is upper semi-continuous, it follows from Cor. 1 that \(f\) is \(\mu\) -integrable for every measure \(\mu\) on \(\mathrm{X}\) . Moreover, \(\mu(f) = \mu(h) - \mu(h - f)\) and \(h - f\) is a lower semi-continuous function \(\geqslant 0\) . Since the mapping \(\mu \mapsto \mu(h - f)\) is lower semi-continuous on \(\mathcal{M}_{+}(\mathrm{X})\) for the vague topology (§1, No.~1, Prop. 4), this proves the corollary.

THEOREM 3. --- For a numerical function \(f \geqslant 0\) to be integrable, it is necessary and sufficient that, for every \(\varepsilon > 0\) , there exist an upper semi-continuous function \(g \geqslant 0\) , with finite values and compact support, and a lower semi-continuous integrable function \(h\) , such that \(g \leqslant f \leqslant h\) and \(\int (h - g) d|\mu| \leqslant \varepsilon\) .

The condition is sufficient by a general criterion for integrability (§3, No.~4, Prop. 8), Prop. 5 and its Corollary 1. Let us show that the condition is necessary. If \(f \geqslant 0\) is integrable then, for every \(\varepsilon > 0\) , there exists a function \(u \geqslant 0\) , continuous and with compact support, such that \(\mathrm{N}_1(f - u) \leqslant \varepsilon / 4\) . By the definition of \(\mathrm{N}_1\) , this implies that there exists a lower semicontinuous function \(v \geqslant 0\) such that \(|\mu|^*(v) \leqslant \varepsilon / 2\) and \(|f - u| \leqslant v\) . Thus, \(-v(x) \leqslant f(x) - u(x) \leqslant v(x)\) for all \(x \in X\) , and since \(u(x)\) is everywhere finite, it follows that \((u(x) - v(x))^+ \leqslant f(x) \leqslant u(x) + v(x)\) for all \(x \in X\) . The functions \(g = (u - v)^+\) and \(h = u + v\) meet the requirements.

COROLLARY. --- For every integrable (resp. integrable and \(\geqslant0\) ) numerical function f, there exist an increasing sequence ( \(g_{n}\) ) of upper semi-continuous functions that are integrable (resp. integrable, with finite values \(\geqslant0\) , and with compact support), and a decreasing sequence ( \(h_{n}\) ) of lower semi-continuous integrable functions, such that:

\(1^{\circ} g_{n}(x) \leqslant f(x) \leqslant h_{n}(x)\) for all \(x \in X\) and every integer \(n\) ;

\(2^{\circ} f(x)\) is equal almost everywhere to the lower envelope h of the sequence \((h_{n})\) and to the upper envelope g of the sequence \((g_{n})\) ;

\[
3 ^ {\circ} \int f d \mu = \lim _ {n \rightarrow \infty} \int g _ {n} d \mu = \lim _ {n \rightarrow \infty} \int h _ {n} d \mu .
\]

Suppose first that \(f \geqslant 0\) . By Th. 3 there exist, for every n, a lower semi-continuous integrable function \(v_{n}\) , and an upper semi-continuous function \(u_{n} \geqslant 0\) with finite values and compact support, such that

\[
u _ {n} \leqslant f \leqslant v _ {n} \quad \text { and } \quad \int (v _ {n} - u _ {n}) d | \mu | \leqslant 1 / n;
\]

setting \(g_{n} = \sup(u_{1}, u_{2}, \ldots, u_{n})\) and \(h_{n} = \inf(v_{1}, v_{2}, \ldots, v_{n})\) , the sequences \((g_{n})\) and \((h_{n})\) meet the requirements. For, since \(g \leqslant f\) , g is integrable by Prop. 4 of No.~3, and since

\[
\int (f - g _ {n}) d | \mu | \leqslant \int (v _ {n} - u _ {n}) d | \mu | \leqslant 1 / n
\]

we have

\[
\int (f - g) d | \mu | = \lim _ {n \rightarrow \infty} \int (f - g _ {n}) d | \mu | = 0
\]

(No.~3, Prop. 4), which proves that \(f\) and \(g\) are equivalent. One argues similarly for the sequence \((h_n)\) .

If \(f\) is not positive, we can apply the foregoing to \(f^{+}\) and \(f^{-}\) , thus there are two increasing sequences \((g_n')\) , \((g_n'')\) of upper semi-continuous integrable functions, and two decreasing sequences \((h_n')\) , \((h_n'')\) of lower semi-continuous integrable functions, such that: \(1^\circ g_n' \leqslant f^{+} \leqslant h_n'\) , \(g_n'' \leqslant -f^{-} \leqslant h_n''\) ; \(2^\circ f^{+}\) (resp. \(-f^{-}\) ) is equal almost everywhere to the upper envelope of the \(g_n'\) and to the lower envelope of the \(h_n'\) (resp. to the upper envelope of the \(g_n''\) and to the lower envelope of the \(h_n''\) ); and \(3^\circ\) :

\[
\begin{array}{r} \int f ^ {+} d \mu = \lim _ {n \to \infty} \int g _ {n} ^ {\prime} d \mu = \lim _ {n \to \infty} \int h _ {n} ^ {\prime} d \mu , \\ - \int f ^ {-} d \mu = \lim _ {n \to \infty} \int g _ {n} ^ {\prime \prime} d \mu = \lim _ {n \to \infty} \int h _ {n} ^ {\prime \prime} d \mu . \end{array}
\]

Moreover, we can suppose that the \(g_{n}^{\prime}\) and the \(h_{n}^{\prime\prime}\) are everywhere finite; it is then clear that the sequences \(g_{n}=g_{n}^{\prime}+g_{n}^{\prime\prime}\) and \(h_{n}=h_{n}^{\prime}+h_{n}^{\prime\prime}\) meet the requirements.

by Th. 2 of No.~3; thus, the integral \(\int_{-\infty}^{+\infty} \mathbf{f}(x) dx\) is absolutely convergent (FRV, II, §2, No.~3). Moreover, \(\int \mathbf{f} d\mu = \int_{-\infty}^{+\infty} \mathbf{f}(x) dx\) by Th. 2 of No.~3. Conversely, suppose that \(\int_{-\infty}^{+\infty} \mathbf{f}(x) dx\) is absolutely convergent; again, by Th. 2 of No.~3, \(\int \mathbf{f} d\mu = \int_{-\infty}^{+\infty} \mathbf{f}(x) dx\) . Note that if the integral \(\int_{-\infty}^{+\infty} \mathbf{f}(x) dx\) is convergent but not absolutely convergent; then \(\mathbf{f}\) is not integrable with respect to Lebesgue measure.

\begin{enumerate}
\def\labelenumi{\arabic{enumi})}
\setcounter{enumi}{1}
\tightlist
\item
  Applied to Lebesgue measure and to regulated functions, Prop. 3 of No.~3 yields anew the theorem on passage to the limit for integrals of regulated functions on a compact interval (FRV, II, §3, No.~1, Prop. 1); for sequences (or for filters with a countable base) of regulated functions, Th. 2 of No.~3 greatly improves on this proposition since, for uniformly bounded regulated functions on a compact interval, it substitutes pointwise convergence for uniform convergence (cf.~§5, No.~4, Th. 2). However, as regards the passage to the limit for absolutely convergent integrals of regulated functions on a noncompact interval, we note that the conditions of Th. 2 of No.~3 imply that the integrals under consideration are uniformly convergent (in the sense defined in FRV, II, §3, No.~2), and thus do not improve on the conditions for convergence given in Book IV (loc. cit.) except as concerns the convergence of the functions \(\mathbf{f}_{\alpha}\) on every compact interval. Finally, the conditions for passing to the limit given for non absolutely convergent integrals of regulated functions remain outside the theory developed in this chapter.
\end{enumerate}

